\section{Level Order Traversal (BFS)}
\label{sec:trees-level-order}

\problemstatement{Given the root of a binary tree, return its level order
traversal: the node values grouped level by level, from the root's level
downward, left to right within each level.}

\begin{patternbox}
\emph{``Level by level,'' ``breadth-first,'' ``top-to-bottom''} are the
signature of a \textbf{queue}-based traversal. Whenever a tree problem
asks for anything organised by \emph{depth} -- level order, left/right
view, bottom/top view, maximum width, minimum time to burn -- a BFS with a
queue is almost always the backbone.
\end{patternbox}

\subsection*{Understanding the problem carefully}

Depth-first recursion dives all the way down one branch before touching
the next, so it visits nodes out of level order. To visit nodes in order
of increasing depth we need the opposite discipline: fully process one
level before starting the next. A \textbf{FIFO queue} does exactly this --
if we always enqueue a node's children when we dequeue it, then all of
level $k$ is enqueued before any node of level $k+1$, so the queue hands
nodes back in level order automatically.

\begin{ideabox}[Freeze the Level Size Before Draining It]
When you begin processing a level, the queue holds \emph{exactly} the
nodes of that level and nothing else. Record its size $s$ first, then
dequeue precisely $s$ nodes: those $s$ are the current level, and the
children you enqueue while draining them form the next level, cleanly
separated. This ``read the size, then loop that many times'' trick is what
lets a single queue produce grouped-by-level output.
\end{ideabox}

\subsection*{Algorithm}

\begin{enumerate}
  \item If root is null, return empty.
  \item Push root into a queue.
  \item While the queue is non-empty:
  \begin{enumerate}
    \item Let $s$ be the current queue size; start a new level list.
    \item Repeat $s$ times: pop the front node, append its value to the
          level list, and push its non-null children.
    \item Append the level list to the answer.
  \end{enumerate}
\end{enumerate}

\subsection*{Dry run}

\dryrun{Tree: root $1$; children $2$ (left) and $3$ (right); $2$'s
children $4,5$.}

\begin{itemize}
  \item Queue $=[1]$, size $1$: pop $1$, push $2,3$. Level $=[1]$.
  \item Queue $=[2,3]$, size $2$: pop $2$ (push $4,5$), pop $3$. Level
        $=[2,3]$.
  \item Queue $=[4,5]$, size $2$: pop $4$, pop $5$. Level $=[4,5]$.
  \item Result: $[[1],[2,3],[4,5]]$.
\end{itemize}

\subsection*{C++ implementation}

\begin{lstlisting}
#include <bits/stdc++.h>
using namespace std;

struct TreeNode {
    int val;
    TreeNode* left;
    TreeNode* right;
    TreeNode(int v) : val(v), left(nullptr), right(nullptr) {}
};

vector<vector<int>> levelOrder(TreeNode* root) {
    vector<vector<int>> ans;
    if (root == nullptr) return ans;

    queue<TreeNode*> q;
    q.push(root);
    while (!q.empty()) {
        int s = q.size();
        vector<int> level;
        for (int i = 0; i < s; i++) {
            TreeNode* node = q.front();
            q.pop();
            level.push_back(node->val);
            if (node->left)  q.push(node->left);
            if (node->right) q.push(node->right);
        }
        ans.push_back(level);
    }
    return ans;
}

int main() {
    TreeNode* root = new TreeNode(1);
    root->left = new TreeNode(2);
    root->right = new TreeNode(3);
    root->left->left = new TreeNode(4);
    root->left->right = new TreeNode(5);

    for (auto& level : levelOrder(root)) {
        for (int x : level) cout << x << " ";
        cout << "\n";
    }
    return 0;
}
\end{lstlisting}

\begin{complexitybox}
\textbf{Time Complexity.} $O(n)$ -- every node is enqueued and dequeued
exactly once. \textbf{Space Complexity.} $O(n)$ for the queue, which in
the worst case (a full last level) holds up to $n/2$ nodes at once.
\end{complexitybox}

\begin{takeawaybox}
The queue plus the ``freeze the level size'' loop is the single reusable
engine behind every by-depth view in this chapter. Get it fluent here;
top view, bottom view, right view and maximum width all reuse it verbatim
with a small tweak to what gets recorded.
\end{takeawaybox}
